% Options for packages loaded elsewhere
\PassOptionsToPackage{unicode}{hyperref}
\PassOptionsToPackage{hyphens}{url}
\documentclass[
  ignorenonframetext,
]{beamer}
\newif\ifbibliography
\usepackage{pgfpages}
\setbeamertemplate{caption}[numbered]
\setbeamertemplate{caption label separator}{: }
\setbeamercolor{caption name}{fg=normal text.fg}
\beamertemplatenavigationsymbolsempty
% remove section numbering
\setbeamertemplate{part page}{
  \centering
  \begin{beamercolorbox}[sep=16pt,center]{part title}
    \usebeamerfont{part title}\insertpart\par
  \end{beamercolorbox}
}
\setbeamertemplate{section page}{
  \centering
  \begin{beamercolorbox}[sep=12pt,center]{section title}
    \usebeamerfont{section title}\insertsection\par
  \end{beamercolorbox}
}
\setbeamertemplate{subsection page}{
  \centering
  \begin{beamercolorbox}[sep=8pt,center]{subsection title}
    \usebeamerfont{subsection title}\insertsubsection\par
  \end{beamercolorbox}
}
% Prevent slide breaks in the middle of a paragraph
\widowpenalties 1 10000
\raggedbottom
\AtBeginPart{
  \frame{\partpage}
}
\AtBeginSection{
  \ifbibliography
  \else
    \frame{\sectionpage}
  \fi
}
\AtBeginSubsection{
  \frame{\subsectionpage}
}
\usepackage{iftex}
\ifPDFTeX
  \usepackage[T1]{fontenc}
  \usepackage[utf8]{inputenc}
  \usepackage{textcomp} % provide euro and other symbols
\else % if luatex or xetex
  \usepackage{unicode-math} % this also loads fontspec
  \defaultfontfeatures{Scale=MatchLowercase}
  \defaultfontfeatures[\rmfamily]{Ligatures=TeX,Scale=1}
\fi
\usepackage{lmodern}
\ifPDFTeX\else
  % xetex/luatex font selection
\fi
% Use upquote if available, for straight quotes in verbatim environments
\IfFileExists{upquote.sty}{\usepackage{upquote}}{}
\IfFileExists{microtype.sty}{% use microtype if available
  \usepackage[]{microtype}
  \UseMicrotypeSet[protrusion]{basicmath} % disable protrusion for tt fonts
}{}
\makeatletter
\@ifundefined{KOMAClassName}{% if non-KOMA class
  \IfFileExists{parskip.sty}{%
    \usepackage{parskip}
  }{% else
    \setlength{\parindent}{0pt}
    \setlength{\parskip}{6pt plus 2pt minus 1pt}}
}{% if KOMA class
  \KOMAoptions{parskip=half}}
\makeatother
\usepackage{longtable,booktabs,array}
\newcounter{none} % for unnumbered tables
\usepackage{calc} % for calculating minipage widths
\usepackage{caption}
% Make caption package work with longtable
\makeatletter
\def\fnum@table{\tablename~\thetable}
\makeatother
\usepackage{graphicx}
\makeatletter
\newsavebox\pandoc@box
\newcommand*\pandocbounded[1]{% scales image to fit in text height/width
  \sbox\pandoc@box{#1}%
  \Gscale@div\@tempa{\textheight}{\dimexpr\ht\pandoc@box+\dp\pandoc@box\relax}%
  \Gscale@div\@tempb{\linewidth}{\wd\pandoc@box}%
  \ifdim\@tempb\p@<\@tempa\p@\let\@tempa\@tempb\fi% select the smaller of both
  \ifdim\@tempa\p@<\p@\scalebox{\@tempa}{\usebox\pandoc@box}%
  \else\usebox{\pandoc@box}%
  \fi%
}
% Set default figure placement to htbp
\def\fps@figure{htbp}
\makeatother
\setlength{\emergencystretch}{3em} % prevent overfull lines
\providecommand{\tightlist}{%
  \setlength{\itemsep}{0pt}\setlength{\parskip}{0pt}}
\usepackage{bookmark}
\IfFileExists{xurl.sty}{\usepackage{xurl}}{} % add URL line breaks if available
\urlstyle{same}
\hypersetup{
  hidelinks,
  pdfcreator={LaTeX via pandoc}}

\author{\texorpdfstring{}{}}
\date{}

\begin{document}

\section{Meaning Spaces as Hilbert Spaces}\label{sec:quantum-semantics}

\begin{frame}[fragile]{\(P(c\mid\rho) = \text{Tr}(E_c\rho)\)}
\protect\phantomsection\label{pcmidrho-texttre_crho}
To connect TQNNs and ZX circuits to \emph{distributional semantics}, we
reinterpret the amplitudes and correlations in these topological
diagrams as semantic quantities. Recent work on quantum semantic
communication supplies the necessary bridge, modeling meaning spaces as
Hilbert spaces at nodes of an interaction graph connected by completely
positive trace-preserving (CPTP) channels along edges
{[}@thomas2025quantum{]}:

\begin{quote}
``Multi-agent semantic networks are modeled as quantum sheaves, where
agents' meaning spaces are Hilbert spaces connected by quantum
channels.'' {[}@thomas2025quantum{]}
\end{quote}

A \emph{quantum semantic sheaf} over a communication graph \(G=(V,E)\)
is a triple \((H,F,\rho)\) where each vertex \(v\) carries a
finite-dimensional semantic Hilbert space \(H_v\), edges carry CPTP maps
\(F_e\), and each vertex holds a density operator \(\rho_v\) encoding
its current semantic state. This instantiates a distributional-semantics
picture: meanings are vectors or density operators in high-dimensional
spaces, with co-occurrence and pragmatic context encoded in the
functorial connections across the network.

\textbf{Case assignment as quantum measurement.} The connection to
classical case systems becomes concrete when we model case assignment as
a \emph{quantum measurement} on the semantic state. Define POVM elements
corresponding to cases
\(\{E_{\text{NOM}}, E_{\text{ACC}}, E_{\text{DAT}}, \ldots\}\)
satisfying \(\sum_c E_c = I\), where each \(E_c\) projects onto the
subspace of semantic states consistent with case role \(c\). The
probability of assigning case \(c\) to a noun phrase in semantic state
\(\rho\) is:

\begin{equation}
P(c \mid \rho) = \text{Tr}(E_c \rho)
\label{eq:eq-8-1}
\end{equation}

\begin{figure}
\centering
\pandocbounded{\includegraphics[keepaspectratio,alt={Overlapping POVM elements produce graded case probabilities via quantum interference. (a) NOM/ACC projectors in a crisp case system: non-overlapping probability density yields deterministic assignment per . (b) Graded case assignment in a proto-role system: overlapping POVM elements create an interference pattern in the semantic belief space, realizing the {[}0,1{]}-enrichment of  as physical measurement. Rotation into a different measurement basis (a different quantum reference frame) corresponds to a different alignment system, e.g., ACC \textbackslash to ERG. Generated programmatically from src/visualization/quantum\_plots.plot\_povm\_probabilities().}]{output/figures/quantum_povm_probabilities.png}}
\caption{Overlapping POVM elements produce graded case probabilities via
quantum interference. (a) NOM/ACC projectors in a crisp case system:
non-overlapping probability density yields deterministic assignment per
\autoref{eq:eq-8-1}. (b) Graded case assignment in a proto-role system:
overlapping POVM elements create an interference pattern in the semantic
belief space, realizing the \([0,1]\)-enrichment of
\autoref{sec:enriched-categories} as physical measurement. Rotation into
a different measurement basis (a different quantum reference frame)
corresponds to a different alignment system, e.g., ACC \(\to\) ERG.
Generated programmatically from
\protect\texttt{src/visualization/quantum\_plots.plot\_povm\_probabilities()}.}\label{fig:quantum-povm}
\end{figure}

For crisp case systems (NOM/ACC), the POVM elements are orthogonal
projectors (\(E_c E_{c'} = \delta_{cc'} E_c\)), yielding deterministic
case assignment. For graded proto-roles (Dowty's
{[}-@dowty1991thematic{]} agent/patient continuum), the POVM elements
overlap, yielding probabilistic case assignment---precisely the quantum
generalization of the \([0,1]\)-enrichment from
\autoref{sec:enriched-categories}. The enriched hom-value
\(\mathcal{C}(v, c)\) is identified with \(P(c \mid \rho_v)\), grounding
the abstract enrichment in physical measurement theory. Fluid-S
alignment (\autoref{sec:case-systems}) then corresponds to a
context-dependent POVM: the measurement basis rotates depending on the
volition feature \(\theta\), so the same noun phrase has different case
probabilities depending on whether the agent is construing the action as
volitional or not.

The same scalar-belief dynamics appear in
\autoref{fig:active-inference-belief}
(\autoref{sec:cognitive-integration}): variational free energy separates
competing case frames before the POVM readout in \autoref{eq:eq-8-1} is
applied to semantic density matrices.
\end{frame}

\begin{frame}{Three Correspondences: Wires, Spiders, and Topology}
\protect\phantomsection\label{three-correspondences-wires-spiders-and-topology}
When this sheaf-theoretic semantics is grafted onto the TQNN/ZX
architecture, three structural correspondences emerge:

\begin{enumerate}
\item
  \textbf{Edges as semantic feature channels}: In a TQNN or ZX-diagram,
  each wire carries not merely an abstract qubit but a semantic Hilbert
  space \(H_v\) associated with a context, concept, or agent. Amplitudes
  or density matrices on that wire encode a distribution over semantic
  features---exactly as word vectors encode distributional meaning in
  classical compositional distributional semantics.
\item
  \textbf{Nodes as compositional operations}: Spiders/gates in ZX or
  intertwiners in TQNNs become \emph{semantic composition maps}: they
  take distributed meanings on input wires and produce new distributed
  meanings on output wires, analogous to how DisCoCat composes word
  meanings into phrase/sentence meanings via multilinear maps.
\item
  \textbf{Topological wiring as contextual structure}: The topology of
  the diagram---the way wires and nodes are connected---encodes which
  semantic spaces interact and in what causal/structural pattern. This
  is the semantic analogue of syntactic structure in distributional
  semantics, realized as a topological quantum circuit.
\end{enumerate}

In this reading, a topological quantum flow neural network becomes a
\emph{distributional semantic machine}: a functor that sends a
topological diagram (graph of contexts and interactions) to a family of
Hilbert spaces and maps where vectors/densities represent distributed
meanings and their probabilistic transformations.
\end{frame}

\begin{frame}{Sheaf Cohomology Governs Alignment}
\protect\phantomsection\label{sheaf-cohomology-governs-alignment}
\begin{block}{Sheaf Cohomology and Semantic Alignment}
\protect\phantomsection\label{sheaf-cohomology-and-semantic-alignment}
The sheaf-based framework proves that semantic alignment between agents
is governed by cohomology classes of the quantum semantic sheaf;
contextuality and entanglement act as resources that remove obstructions
to alignment {[}@thomas2025quantum{]}:

\begin{quote}
``We derive semantic channel capacity when sender and receiver share
prior entanglement, proving it strictly exceeds classical capacity. The
quantum advantage grows as channel noise increases---precisely when
semantic communication most benefits over bit-level transmission.''
{[}@thomas2025quantum{]}
\end{quote}

Two further results from this framework are particularly relevant:

\begin{itemize}
\tightlist
\item
  \textbf{Contextuality as a semantic resource}: ``Quantum contextuality
  reduces cohomological obstructions to semantic alignment. Contextual
  correlations act as `pre-shared semantic resolution,' establishing
  contextuality as a resource for semantic communication''
  {[}@thomas2025quantum{]}.
\item
  \textbf{Discord as integrated semantic information}: ``Quantum discord
  equals integrated semantic information, linking quantum correlations
  to irreducible semantic content and connecting our framework to
  integrated information theory'' {[}@thomas2025quantum{]}.
\end{itemize}

These results establish that the topology of the semantic sheaf (and its
cohomology) constrains how probabilistic semantic information can be
transferred; quantum features (entanglement, contextuality, discord)
change these constraints in well-defined ways.
\end{block}

\begin{block}{Topological Circuit as Semantic Sheaf Skeleton}
\protect\phantomsection\label{topological-circuit-as-semantic-sheaf-skeleton}
A TQNN/ZX circuit implementing a quantum communication or computation
protocol is itself a diagram over which one can define a sheaf of
semantic spaces and channels. The underlying graph of the TQNN/ZX
diagram serves as the base graph \(G=(V,E)\) of the semantic sheaf:
vertices inherit meaning spaces \(H_v\), edges inherit CPTP maps
\(F_e\), and the TQFT/ZX functor gives the global linear map
representing the protocol. Distributional semantics as diagram
evaluation then becomes literal: passing an initial semantic state
(distribution over meanings) through the TQNN/ZX diagram yields an
output state whose components encode the posterior semantic
distributions at boundary wires.

ZX rewrite rules, which change the internal topology of the diagram
while preserving its overall semantics as a linear map, correspond to
alternative factorizations of the same semantic
transformation---different internal ``flow architectures'' for realizing
the same semantic map.
\end{block}
\end{frame}

\begin{frame}{Case Assignment as Holographic Measurement}
\protect\phantomsection\label{case-assignment-as-holographic-measurement}
\begin{block}{Case Assignment as Quantum Measurement}
\protect\phantomsection\label{case-assignment-as-quantum-measurement}
The active inference model of case reasoning
(\autoref{sec:cognitive-integration}) acquires a new dimension in this
quantum topological setting. Case assignment---the cognitive process of
determining \emph{who does what to whom}---can be modeled as a quantum
measurement process on a holographic screen, with the following
correspondences:

{\def\LTcaptype{none} % do not increment counter
\begin{longtable}[]{@{}
  >{\raggedright\arraybackslash}p{(\linewidth - 2\tabcolsep) * \real{0.5000}}
  >{\raggedright\arraybackslash}p{(\linewidth - 2\tabcolsep) * \real{0.5000}}@{}}
\toprule\noalign{}
\begin{minipage}[b]{\linewidth}\raggedright
Classical Case Assignment
\end{minipage} & \begin{minipage}[b]{\linewidth}\raggedright
Quantum Topological Model
\end{minipage} \\
\midrule\noalign{}
\endhead
Case role (NOM, ACC, \ldots) & Pointer state selected by QRF \\
Case frame (alignment system) & Quantum reference frame \\
Relational structure of event & Spin-network topology \\
Free-energy minimization & TQFT evaluation of diagram \\
Prediction-error (P600/N400) & Symmetry-breaking on holographic
screen \\
\bottomrule\noalign{}
\end{longtable}
}
\end{block}

\begin{block}{From Predictive Processing to Topological Flow}
\protect\phantomsection\label{from-predictive-processing-to-topological-flow}
In the predictive processing account, a cognitive agent maintains a
generative model that predicts the relational structure of incoming
linguistic material. When this model is realized as a TQNN, prediction
becomes evaluation of the topological diagram; prediction error becomes
the discrepancy between the predicted TQFT evaluation and the observed
data; and belief updating becomes modification of the spin-network's
edge labels (representation labels) and vertex intertwiners.

The topological character of this computation confers significant
advantages for active inference on case structure: topological
invariants are robust to continuous deformation, so the generative
model's predictions are stable under small perturbations of the
input---a desirable property for language understanding in noisy
environments.
\end{block}

\begin{block}{Entanglement Strictly Exceeds Classical Semantic Capacity}
\protect\phantomsection\label{entanglement-strictly-exceeds-classical-semantic-capacity}
The sheaf-theoretic results of Thomas and Chen {[}-@thomas2025quantum{]}
suggest that quantum features provide genuine advantages for semantic
communication---not merely computational speedup, but qualitative
enhancements in semantic alignment. If case-marked relational structure
is communicated between agents via quantum channels, entanglement
provides additional semantic capacity, contextuality removes alignment
obstructions, and discord captures irreducible semantic content. These
are not abstract possibilities but operational consequences of the
mathematical framework developed across this review. Recent work by
Krawchuk et al. {[}-@krawchuk2025paramqnlp{]} demonstrates this
concretely: DisCoCirc string diagrams that represent discourse-level
semantics (including case role assignments across sentences) can be
\emph{automatically} compiled into parameterized quantum circuits,
closing the loop from linguistic case structure through categorical
formalism to quantum computation.
\end{block}
\end{frame}

\end{document}
